% Generated draft: requires mathematical and source review.
\documentclass[11pt]{article}
\usepackage[margin=25mm]{geometry}
\usepackage{fontspec}
\setmainfont{Latin Modern Roman}
\usepackage{amsmath,amssymb,mathtools,unicode-math}
\setmathfont{Latin Modern Math}
\usepackage{enumitem,needspace,xurl,hyperref}
\setlength{\parindent}{0pt}
\setlength{\parskip}{5pt}
\setlength{\emergencystretch}{3em}
\newcommand{\N}{\mathbb{N}}
\newcommand{\Z}{\mathbb{Z}}
\newcommand{\Q}{\mathbb{Q}}
\newcommand{\R}{\mathbb{R}}
\newcommand{\C}{\mathbb{C}}
\newcommand{\F}{\mathbb{F}}
\newcommand{\A}{\mathbb{A}}
\newcommand{\PP}{\mathbb{P}}
\newcommand{\E}{\mathbb{E}}
\newcommand{\eps}{\varepsilon}
\newcommand{\dd}{\,\mathrm{d}}
\begin{document}
\section*{UnsolvedMath ID: 700659}
\textbf{Source problem code:} catalogue entry
\subsection{Definitions and mathematical statement}
Let \(R\) and \(S\) be associative rings with identity, and assume that \(2\cdot 1_S\) is invertible in \(S\). For positive matrix sizes \(n,m\), put \(V_R=R^n_R\) and \(W_S=S^m_S\). Write \(GL_n(R)=\operatorname{Aut}_R(V)\), and use the usual elementary subgroup
\[
E_n(R)=\left\langle I_n+r e_{ij}:r\in R,\ 1\leq i,j\leq n,\ i\neq j\right\rangle,
\]
where \(e_{ij}\) is the matrix unit; define \(E_m(S)\) similarly. The source gives no further restrictions on \(n\) or \(m\).

The source asks to prove that every \emph{standard homomorphism}
\[
\Lambda:E_n(R)\longrightarrow E_m(S)
\]
originates from a collineation and a correlation: there should exist maps described in the source as semilinear isomorphisms
\[
g:V_R\longrightarrow W_S,
\qquad
h:V_R\longrightarrow {}_S\!W'_{S^0},
\]
and an object \(f\) such that, for every \(x\in E_n(R)\),
\[
\Lambda(x)=(g^{-1}xg)f+(h^{-1}xh)(1-f).
\]
Here \(g\) is called a collineation and \(h\) a correlation.

\textit{Unresolved qualifications.} This is a transcription of the proposed assertion, not a fully specified proposition. The supplied record does not define ``standard homomorphism,'' specify the scalar maps underlying semilinearity, define \(W'\) or \(S^0\), or state the ambient algebra and required properties of \(f\). In particular, no idempotence or centrality condition on \(f\) is supplied. The bimodule notation in the target of \(h\), the induced action on that target, and the composition convention needed to interpret the displayed conjugations are also unspecified. Reading \(S^0\) as an opposite ring and \(W'\) as a dual module would require confirmation. These omissions prevent a precise, unqualified interpretation of the asserted formula.
\subsection{Short English statement}
Petechuk's problem asks whether every standard homomorphism between the stated elementary matrix groups over unital associative rings, with \(2\) invertible in the target ring, has the displayed decomposition using a collineation and a correlation. Essential terminology and the parameter \(f\) are undefined in the supplied record, so the intended assertion requires clarification. The registry labels the problem open and separately reports an unverified AI counterexample claim.
\subsection{Sources}
\begin{itemize}
\item Local UnsolvedMath registry, permanent ID \(700659\), source code \texttt{catalogue entry}, title ``Kourovka Notebook Problem 10.44.'' The supplied record attributes the problem to V. M. Petechuk, proposed in 1986.
\item Supplied source attribution: \emph{The Kourovka Notebook (Unsolved Problems in Group Theory)}, 10th issue (1986), Problem 10.44; 21st issue edition (2026), \texttt{arXiv:1401.0300v48}. The registry's collection description names E. I. Khukhro and V. D. Mazurov as editors.
\item The supplied background and status-review notes report: an unverified AI claim that a counterexample was constructed, attributed to Achyuth Jayadevan, \emph{Preprint}, 2026, \url{https://zenodo.org/records/22847412}. The status-review notes say that the problem is marked with \(\#\) for this unverified claim in the October 2026 21st issue. This is supplied attribution, not verification of a counterexample.
\item The supplied status review is dated \texttt{2026-10-16}, records canonical status \texttt{open}, and classifies its evidence as \texttt{unrefereed\_claim}. This assessment is supplied dated curation, not independently verified. No cited source was independently consulted for this draft.
\end{itemize}
\end{document}
